\bibitem[1]{matrix-ref-1} Samson Abramsky, Radha Jagadeesan, and Pasquale Malacaria. Full abstraction for PCF.
Information and Computation, 163(2):409–470, December 2000. doi:10.1006/inco.2000.
2930.
\bibitem[2]{matrix-ref-2} Patrick Baillot, Vincent Danos, Thomas Ehrhard, and Laurent Regnier. Timeless games.
In Computer Science Logic, pages 56–77. Springer Berlin Heidelberg, 1998. doi:10.1007/
bfb0028007.
\bibitem[3]{matrix-ref-3} François Bergeron, Gilbert Labelle, and Pierre Leroux. Combinatorial Species and Tree-like
Structures. Cambridge University Press, November 1997. doi:10.1017/cbo9781107325913.
\bibitem[4]{matrix-ref-4} Gérard Berry. Stable models of typed \ensuremath{\lambda}-calculi, pages 72–89. Springer Berlin Heidelberg, 1978.
doi:10.1007/3-540-08860-1\_7.
\bibitem[5]{matrix-ref-5} Pierre Boudes. Projecting games on hypercoherences. In International Colloquium on Automata,
Languages and Programming, pages 257–268. Springer Berlin Heidelberg, 2004. doi:10.1007/
978-3-540-27836-8\_24.
\bibitem[6]{matrix-ref-6} Jean Bénabou. Introduction to bicategories, pages 1–77. Springer Berlin Heidelberg, 1967.
doi:10.1007/bfb0074299.
\bibitem[7]{matrix-ref-7} Ana C. Calderon and Guy McCusker. Understanding game semantics through coherence
spaces. Electronic Notes in Theoretical Computer Science, 265:231–244, September 2010.
doi:10.1016/j.entcs.2010.08.014.
\bibitem[8]{matrix-ref-8} Pierre Clairambault and Marc de Visme. Full abstraction for the quantum lambda-calculus.
Proceedings of the ACM on Programming Languages, 4(POPL):1–28, December 2020. doi:
10.1145/3371131.
\bibitem[9]{matrix-ref-9} Pierre Clairambault and Simon Forest. The cartesian closed bicategory of thin spans of
groupoids. In 2023 38th Annual ACM/IEEE Symposium on Logic in Computer Science
(LICS), pages 1–13. IEEE, June 2023. doi:10.1109/lics56636.2023.10175754.
\bibitem[10]{matrix-ref-10} Pierre Clairambault and Simon Forest. An analysis of symmetry in quantitative semantics. In
Proceedings of the 39th Annual ACM/IEEE Symposium on Logic in Computer Science, LICS
’24, pages 1–13. ACM, July 2024. doi:10.1145/3661814.3662092.
\bibitem[11]{matrix-ref-11} Pierre Clairambault, Marc De Visme, and Glynn Winskel. Game semantics for quantum
programming. Proceedings of the ACM on Programming Languages, 3(POPL):1–29, January
2019. doi:10.1145/3290345.
\bibitem[12]{matrix-ref-12} J.A.B. Cockett, M. Hasegawa, and R.A.G. Seely. Coherence of the double involution on
\ensuremath{*}-autonomous categories. Theory and Applications of Categories, 17:17–29, 2006. URL:
http://eudml.org/doc/127888.
\bibitem[13]{matrix-ref-13} Thomas Ehrhard and Laurent Regnier. Uniformity and the taylor expansion of ordinary
lambda-terms. Theoretical Computer Science, 403(2-3):347–372, August 2008. doi:10.1016/
j.tcs.2008.06.001.
\bibitem[14]{matrix-ref-14} M. Fiore, N. Gambino, and M. Hyland. Monoidal bicategories, differential linear logic, and
analytic functors, 2024. doi:10.48550/arXiv.2405.05774.
\bibitem[15]{matrix-ref-15} M. Fiore, N. Gambino, M. Hyland, and G. Winskel. The cartesian closed bicategory of
generalised species of structures. Journal of the London Mathematical Society, 77(1):203–220,
October 2007. doi:10.1112/jlms/jdm096.
\bibitem[16]{matrix-ref-16} Marcelo Fiore, Zeinab Galal, and Hugo Paquet. A combinatorial approach to higher-order
structure for polynomial functors. In FSCD. Schloss Dagstuhl – Leibniz-Zentrum für Informatik,
2022. doi:10.4230/LIPIcs.FSCD.2022.31.
\bibitem[17]{matrix-ref-17} Marcelo Fiore, Zeinab Galal, and Hugo Paquet. Stabilized profunctors and stable species
of structures. Logical Methods in Computer Science, Volume 20, Issue 1, February 2024.
doi:10.46298/lmcs-20(1:17)2024.
\bibitem[18]{matrix-ref-18} Jean-Yves Girard. Linear logic. Theoretical Computer Science, 50(1):1–101, 1987. doi:
10.1016/0304-3975(87)90045-4.
\bibitem[19]{matrix-ref-19} Jean-Yves Girard. Normal functors, power series and \ensuremath{\lambda}-calculus. Annals of Pure and Applied
Logic, 37(2):129–177, February 1988. doi:10.1016/0168-0072(88)90025-5.
\bibitem[20]{matrix-ref-20} J.M.E. Hyland and C.-H.L. Ong. On full abstraction for PCF: I, II, and III. Information and
Computation, 163(2):285–408, December 2000. doi:10.1006/inco.2000.2917.
\bibitem[21]{matrix-ref-21} Martin Hyland. Abstract interpretation of proofs: Classical propositional calculus. In
Computer Science Logic (CSL 2004), pages 6–21. Springer Berlin Heidelberg, 2004. doi:
10.1007/978-3-540-30124-0\_2.
\bibitem[22]{matrix-ref-22} Martin Hyland and Andrea Schalk. Glueing and orthogonality for models of linear logic. The-
oretical Computer Science, 294(1-2):183–231, February 2003. doi:10.1016/s0304-3975(01)
00241-9.
\bibitem[23]{matrix-ref-23} André Joyal. Une théorie combinatoire des séries formelles. Advances in Mathematics,
42(1):1–82, October 1981. doi:10.1016/0001-8708(81)90052-9.
\bibitem[24]{matrix-ref-24} Jim Laird, Giulio Manzonetto, Guy McCusker, and Michele Pagani. Weighted relational
models of typed lambda-calculi. In 2013 28th Annual ACM/IEEE Symposium on Logic in
Computer Science. IEEE, June 2013. doi:10.1109/lics.2013.36.
\bibitem[25]{matrix-ref-25} François Lamarche. Quantitative domains and infinitary algebras. Theoretical Computer
Science, 94(1):37–62, March 1992. doi:10.1016/0304-3975(92)90323-8.
\bibitem[26]{matrix-ref-26} Paul-André Melliès. Asynchronous games 1: Uniformity by group invariance, 2003. URL:
http://www.pps.univ-paris-diderot.fr/~mellies/papers/ag1.pdf.
\bibitem[27]{matrix-ref-27} Paul-André Melliès. Sequential algorithms and strongly stable functions. Theoretical Computer
Science, 343(1–2):237–281, October 2005. doi:10.1016/j.tcs.2005.05.015.
\bibitem[28]{matrix-ref-28} Paul-André Melliès. Categorical semantics of linear logic. In Interactive models of computation
and program behaviour, panoramas et synthèses, volume 27, pages 1–196. Société Mathématique
de France, 2009.
\bibitem[29]{matrix-ref-29} Federico Olimpieri. Intersection type distributors. In 2021 36th Annual ACM/IEEE Symposium
on Logic in Computer Science (LICS), pages 1–15. IEEE, June 2021. doi:10.1109/lics52264.
2021.9470617.
\bibitem[30]{matrix-ref-30} Christian Retoré. A semantic characterisation of the correctness of a proof net. Math-
ematical Structures in Computer Science, 7(5):445–452, October 1997. doi:10.1017/
s096012959700234x.
\bibitem[31]{matrix-ref-31} Andrea Schalk. What is a categorical model for linear logic? Available on the author’s web
page: https://www.cs.man.ac.uk/~schalk/notes/index.html, 2004.
\bibitem[32]{matrix-ref-32} Paul Taylor. Quantitative domains, groupoids and linear logic. In David H. Pitt, David E.
Rydeheard, Peter Dybjer, Andrew M. Pitts, and Axel Poigné, editors, Category Theory
and Computer Science, pages 155–181, Berlin, Heidelberg, 1989. Springer Berlin Heidelberg.
doi:10.1007/BFB0018351.
\bibitem[33]{matrix-ref-33} Takeshi Tsukada, Kazuyuki Asada, and Kengo Hirata. Full definability in a profunctorial
model, 2026. doi:10.48550/arXiv.2604.26829.
\bibitem[34]{matrix-ref-34} Takeshi Tsukada, Kazuyuki Asada, and C.-H. Luke Ong. Generalised species of rigid resource
terms. In 2017 32nd Annual ACM/IEEE Symposium on Logic in Computer Science (LICS).
IEEE, June 2017. doi:10.1109/lics.2017.8005093.
\bibitem[35]{matrix-ref-35} Takeshi Tsukada, Kazuyuki Asada, and C.-H. Luke Ong. Species, profunctors and taylor
expansion weighted by SMCC. In Proceedings of the 33rd Annual ACM/IEEE Symposium on
Logic in Computer Science. ACM, July 2018. doi:10.1145/3209108.3209157.
\bibitem[36]{matrix-ref-36} Takeshi Tsukada and C.-H. Luke Ong. Nondeterminism in game semantics via sheaves. In
2015 30th Annual ACM/IEEE Symposium on Logic in Computer Science. IEEE, July 2015.
doi:10.1109/lics.2015.30.
\bibitem[37]{matrix-ref-37} Takeshi Tsukada and C.-H. Luke Ong. Plays as resource terms via non-idempotent intersection
types. In Proceedings of the 31st Annual ACM/IEEE Symposium on Logic in Computer
Science. ACM, July 2016. doi:10.1145/2933575.2934553.
